\section{Introduction}
The scheduling of maintenance operations for large-scale transmission networks constitutes a critical challenge in modern energy systems. As grid operators face the dual pressure of aging infrastructure and the integration of intermittent renewable energy sources, the necessity to balance operational safety, financial cost, and environmental impact becomes paramount. While preventive maintenance is essential for grid reliability, the act of temporarily removing a power line from service weakens the network's ability to respond to unexpected contingencies such as extreme weather events or cascading equipment failures. This operational tension has been amplified by the energy transition, as the increasing penetration of intermittent renewable generation creates new dynamics in power flows and demand patterns.

To tackle this problem, RTE (Réseau de Transport d'Électricité), the French transmission system operator, developed a framework that first quantifies the financial risk of scheduling each maintenance task at each possible time period, evaluated across thousands of future demand and generation scenarios. The combinatorially complex problem of using this pre-calculated risk data to construct a feasible annual maintenance schedule was the focus of the ROADEF/EURO Challenge 2020 \citep{roadef2020}, which sought algorithmic approaches capable of managing thousands of maintenance tasks under strict temporal and resource constraints. 

In the context of such competitions, and in combinatorial optimization generally, the objective function is frequently simplified to ensure computational tractability. The ROADEF 2020 challenge quantified solution quality via a convex combination of mean estimated risk and a tail-risk metric (the excess of the 90th percentile). While mathematically convenient for guiding local search heuristics and integer programming solvers, this aggregated metric collapses the multidimensional nature of the decision problem into a single scalar value. Consequently, near optimal solutions do not prevent operational fragility when subjected to the nuanced preferences of a decision maker. The optimization phase, therefore, is best viewed not as the final decision step, but as a generator of a set of feasible candidate schedules, a choice set, from which a final selection must be made.

This realization necessitates a transition from combinatorial optimization to Multi-Criteria Decision Making (MCDM). Within the taxonomy of MCDM problems described by \citet{roy1996multicriteria}, this study addresses the choice problem%($P.\alpha$)
: the selection of the single most robust schedule from a finite set of high-quality alternatives generated by the top-performing algorithms from the challenge, including the approach by \citet{buljubasic2020} and the metaheuristics developed by \citet{parreno2020} and \citet{langiu2020}.

To aggregate the heterogeneous criteria defining these schedules ranging from grid reliability to socio-environmental externalities, the Ordered Weighted Averaging (OWA) operator, introduced by \citet{yager1988}, serves as a foundational tool. Unlike the classical weighted mean, which assumes constant trade-offs between criteria, the OWA operator enables the modeling of complex attitudinal characters, such as risk aversion or optimism, by reordering arguments prior to aggregation \cite{emrouznejad2014ordered}. The capacity of OWA to represent varying degrees of compensation among attributes has been empirically validated as a realistic proxy for human preference structures in ranking tasks \cite{reimann2017how}. Furthermore, given the inherent imprecision in quantifying spatial impacts (e.g., proximity to protected nature reserves), fuzzy set theory provides a rigorous formalism for mapping raw operational data into semantic criteria suitable for aggregation \cite{zadeh1965fuzzy, cinelli2020classification}.

However, standard OWA aggregation treats criteria symmetrically based solely on their magnitude, ignoring the structural priority often required in safety-critical domains where ``Safety'' must strictly precede ``Cost'' or ``Strategy''. To address this, the Weighted OWA (WOWA) operator \cite{yager2008prioritized} assigns an importance value to each criterion, constraining the compensatory influence of lower-priority attributes. 

In practical settings, the precise numerical weights representing the decision maker's attitude are rarely known with certainty. A solution that appears optimal under a specific weight vector may be unstable, losing its dominance with minor perturbations in the decision maker's preferences.

This paper bridges the gap between large-scale optimization and robust decision analysis. A novel decision support framework is proposed that transforms the outputs of the ROADEF/EURO 2020 solvers into a multi-criteria decision matrix. The primary contribution, however, is methodological and addresses a precise obstacle in the practical use of WOWA for robust choice. Under WOWA, each alternative induces its own cumulative-importance partition and therefore its own effective weight vector. Consequently, pairwise indifference is not naturally represented in a shared weight space, and the geometric stability analysis available for additive models cannot be transferred directly. For a fixed finite choice set, a transformation is derived that refines those alternative-specific partitions into a unified grid and maps WOWA into an equivalent high-dimensional OWA on a common simplex. This exact equivalence restores a shared polyhedral geometry. Robust selection can then be formulated as an exact Linear Programming problem rather than relying on family-dependent parameter sweeps or sampled sensitivity analyses over quantifiers. The proposed model demonstrates that solutions minimizing the original linear objective function from the competition may be rejected when the priority of grid safety is rigorously enforced, highlighting the necessity of rigorous decision analysis in critical infrastructure planning.
